\chapter*{\textarabic{ملخص}}
\addcontentsline{toc}{chapter}{arabic Abstract}

\begingroup
\setstretch{1.12}

\begin{otherlanguage}{arabic}
يقدم هذا التقرير مشروع نهاية الدراسة المنجز داخل شركة \textenglish{\textbf{Scientiae}} والمتعلق بتصميم وتنفيذ منصة \textenglish{\textit{Graph-RAG}} متعددة الوسائط لاستكشاف الرسم البياني المعرفي والاستعلام عنه مع توليد إجابات موثقة. يجمع النظام بين الزحف الذاتي على المصادر السياسية، والبحث الدلالي الكثيف فوق المجموعة النصية المجمعة، والاجتياز المنظم لرسم بياني معرفي في \textenglish{\textbf{Neo4j}} يضم كيانات وعلاقات سياسية مستخرجة آليا، وتوليف إجابات مؤسسة على الأدلة بمساعدة نموذج \textenglish{\textbf{Mistral}}. يهدف المشروع إلى تحويل المعلومات السياسية غير المهيكلة إلى طبقة معرفية قابلة للتتبع والاستكشاف والتحليل.

يعالج المشروع محدوديات الجيل المعزز بالاسترجاع الكلاسيكي عند تطبيقه على المعطيات السياسية ذات البعد العلائقي. يمكن للاسترجاع التقليدي إيجاد مقاطع ذات صلة، لكنه يواجه صعوبة في الأسئلة متعددة القفزات، وإزالة لبس الكيانات، والتفسير الزمني، وتتبع العلاقات بشكل صريح. ومن خلال دمج الاسترجاع الشعاعي مع أدلة منظمة على شكل رسم بياني، تحسن المنصة المقترحة القدرة على الإجابة عن أسئلة تحليلية معقدة مع الحفاظ على إمكانية الرجوع إلى الوثائق المصدرية.

تعتمد المنصة على معمارية تفصل بين جمع المعطيات، والأوركسترا غير المتزامنة، ومعالجة اللغة الطبيعية، والتخزين الوسيط، ثم النشر في \textenglish{\textbf{Neo4j}} وفهرسة المقاطع في \textenglish{\textbf{ZVec}}. وتشمل سلسلة الجمع مراقبة تدفقات \textenglish{\textbf{RSS}} المنتقاة، واكتشاف الأخبار عبر \textenglish{Google News}، وتنزيل وتفكيك تقارير \textenglish{PDF}، مع لوحة إدارة للجدولة والإشراف. وتشمل سلسلة المعالجة التنظيف، واستخراج الكيانات والعلاقات، وحل الإحالات المرجعية بشكل حتمي، وتخزين الوثائق في \textenglish{\textbf{MongoDB}}، ثم تمكين الاستعلام عبر واجهة \textenglish{\textbf{React}} بأنماط تفكير متعددة (\textenglish{rag}، \textenglish{cot}، \textenglish{reflexive}) مع إمكانية تعزيز الاسترجاع ببحث ويب مباشر.

أجري تقييم المنصة على مجموعة من أربعة عشر سؤالا مصنفة (واقعي، علائقي، زمني، مقارن) بمقاييس من نوع \textenglish{\textbf{RAGAS}} بعد تحسين الموجهات. وأظهرت النتائج تحسنا تدريجيا من نمط \textenglish{RAG} إلى نمط التفكير الانعكاسي، مع بقاء استدعاء السياق محورا للتحسين المستقبلي. ويُنظر إلى الرسم البياني المعرفي والمنصة المحيطة به بوصفهما طبقة معرفية قابلة لإعادة الاستعمال في الاستشارة والتحليل السياسي والاستغلال المنظم للمعلومة.
\end{otherlanguage}

\noindent\rule[2pt]{\textwidth}{0.5pt}

\begin{otherlanguage}{arabic}
{\textbf{الكلمات المفتاحية:}} \textenglish{Graph-RAG}، رسم بياني معرفي، \textenglish{Neo4j}، \textenglish{Mistral}، زحف ذاتي، استخبارات سياسية، معالجة اللغة الطبيعية، استخراج العلاقات، استرجاع معزز بالتوليد، أنظمة موزعة.
\end{otherlanguage}

\noindent\rule[2pt]{\textwidth}{0.5pt}
\endgroup
